% ======================================================================
%  Section 4.7.1 --- revised: prose + corrected Table 4.1
%
%  Replaces the current 4.7.1, which has figures and a table but no text.
%
%  UNIT CORRECTIONS APPLIED (see the accompanying message):
%    vpPoseVector stores translation in METRES and rotation as a theta-u
%    vector in RADIANS. The previous table carried the raw console values
%    under "mm" and "deg" headers, making every translation 1000x too
%    small and every rotation 57.3x too small.
%
%  All values are at iteration 600 for all three variants.
%  STILL OPEN, marked %%% OPEN below: whether Variants 2 and 3 genuinely
%  coincide, or whether the two consoles supplied are the same run.
% ======================================================================

\subsection{Nominal case}
\label{sec:nominal}

Before occlusion is introduced, it is necessary to establish what the three
variants do when there is nothing to reject. Two distinct questions are settled
by this run. The first is what the weighted formulations cost when there is
nothing to reject: Remark~\ref{rem:baseline} establishes algebraically that
$\mathbf{D}=\mathbf{I}$ reduces~\eqref{eq:robust-lm-law} to the control law of
Chapter~\ref{ch:herpvs}, but says nothing about $\mathbf{D}\neq\mathbf{I}$ on
uncorrupted measurements, and that is what the nominal run measures. The second
is a matter of attribution: the nominal behaviour of each variant must be known
before any difference observed under occlusion can be ascribed to the treatment
of the occlusion rather than to the weighting altering the control law
independently of it. As will be seen, the three variants do \emph{not} coincide
on a clean scene, which makes this baseline indispensable rather than merely
confirmatory.

\paragraph{Conditions.}
The reference image $\mathbf{I}^{*}$ is acquired from the unoccluded target at
the desired pose $\mathbf{t}^{*}=(0,0,1300)\,\mathrm{mm}$,
$\boldsymbol{\omega}^{*}=(0,0,0)^{\circ}$, at which the target overfills the
image. For the nominal run alone, the images acquired during servoing are also
taken from the unoccluded target, so that the residual arises solely from the
pose error and no measurement is corrupted. Servoing starts from
$\mathbf{t}_{0}=(40,-40,1100)\,\mathrm{mm}$,
$\boldsymbol{\omega}_{0}=(12,-12,6)^{\circ}$, a displacement of approximately
$41$~pixels in the mean and $65$~pixels at worst, with the target covering the
whole image at both poses so that no measurement is lost off-frame. All
remaining quantities --- the Hermite bank and its four scales, the interaction
matrix~\eqref{eq:stacked-L}, the adaptive gain and
damping~\eqref{eq:lambda}--\eqref{eq:mu}, the Tukey constant, the sampling
period and the iteration cap --- are those fixed in
Section~\ref{sec:ablation-protocol} and are identical across the three variants.
That the three runs did begin from the same image is confirmed by the reference
error norm $N_{ei}$, which took the same value $4.04\times10^{8}$ in all three.

\paragraph{Results.}
Figures~\ref{fig:v1}--\ref{fig:v3} show the camera velocities, the positioning
error and the feature error against iteration for the three variants, and
Table~\ref{tab:ch4-nominal} collects the final pose errors. All three variants
position the camera to better than one millimetre and better than one degree,
which is the accuracy expected of a direct scheme at convergence. Within that,
Variant~3 attains the smallest rotational error of the three by a wide margin:
$0.089^{\circ}$, against $0.219^{\circ}$ for the unweighted baseline and
$0.601^{\circ}$ for plain Tukey weighting. Its translational error,
$0.825$~mm, matches Variant~2's and is dominated by the depth component, the
weakest-conditioned direction of a direct scheme.

Two observations require comment, as both bear on how the remaining sections
must be read.

First, the feature error does not decay to zero but settles on a floor:
$4.35\times10^{7}$ for Variant~1, $1.04\times10^{8}$ for Variant~2 and
$6.59\times10^{7}$ for Variant~3, while the commanded velocity falls to between
$2.9\times10^{-5}$ and $3.7\times10^{-4}\,\mathrm{m\,s^{-1}}$ and the pose
error becomes sub-millimetre. The quantity reported is the sum of squares
$\norm{\mathbf{e}}^{2}$ rather than the norm, and the feature has dimension
$264\,000$ --- four Hermite scales over the $220\times300$ interior left by the
ten-pixel filter border --- so the floor corresponds to a root-mean-square
residual of $19.8$ per component, against $1149$ at the first iteration. The
run terminates not because the residual has been minimised but
because the adaptive gain has decayed with it: by~\eqref{eq:lambda} the gain is
proportional to the intensity error, $\lambda=\norm{\mathbf{e}_I}^{2}/
N_{e}^{0.8}$, so it falls from $\lambda\simeq53$ at the first iteration to
$\lambda=0.057$ at iteration~600, a reduction by a factor of $9\times10^{2}$.
Since the commanded velocity is in addition proportional to the residual
itself, $\norm{\mathbf{V}_c}$ decays much faster than the error and had reached
$3.7\times10^{-4}\,\mathrm{m\,s^{-1}}$ while a small misalignment remained.
The minimisation is not at fault: over the same run the feature error fell from
$3.48\times10^{11}$ to $1.04\times10^{8}$, by three and a half orders of
magnitude, the commanded velocity from $0.566$ to
$3.7\times10^{-4}\,\mathrm{m\,s^{-1}}$, and the intensity error from $72.7$ to
$2.40$ grey levels root-mean-square.
It is also worth recording that the damping term becomes inactive very early.
By~\eqref{eq:mu}, $\mu\propto\norm{\mathbf{e}_I}^{10}$, so the same factor of
$9\times10^{2}$ in the intensity error drives $\mu$ from $1.9\times10^{-2}$ to
$2.9\times10^{-17}$: the scheme is a Levenberg--Marquardt method only during
the first few iterations and a Gauss--Newton method thereafter, which is why
the fixed floor added to the damped Hessian in Section~\ref{sec:implementation}
is what actually guarantees invertibility once robust rejection is active.
That misalignment is dominated
by its rotational part: at $f=870$~px the final rotational errors correspond to
image displacements of $3.33$~pixels for Variant~1, $9.13$ for Variant~2 and
$1.35$ for Variant~3, whereas the translational errors, all below a
millimetre at $Z=1.3$~m, correspond to less than a tenth of a pixel.

These displacements account for the floors in order of magnitude but not in
detail. The floors rank $4.35\times10^{7}$, $6.59\times10^{7}$,
$1.04\times10^{8}$ for Variants~1,~3,~2, whereas the misalignments rank
$1.35$, $3.33$, $9.13$~pixels for Variants~3,~1,~2: Variant~3 attains the
smallest misalignment but only the middle floor. The discrepancy is to be
expected, since the residual of a weighted variant is shaped not only by the
misalignment but by which measurements the estimator retains, and the floor
therefore mixes an accuracy term with a coverage term. The feature error is
accordingly reported throughout this chapter as a convergence indicator
only; the pose error of Table~\ref{tab:ch4-nominal} is the accuracy metric.
The
absolute threshold $\norm{\mathbf{e}}^{2}>10^{4}$ inherited from the
implementation lies more than three orders of magnitude below this floor and is
never reached, so convergence is reported throughout this chapter against the
commanded velocity, a run being taken as converged once
$\norm{\mathbf{V}_c}$ falls below $10^{-3}\,\mathrm{m\,s^{-1}}$. The criterion
stated in Section~\ref{sec:ablation-protocol} is to be read accordingly.

Second, the residual normalization is \emph{not} transparent in the nominal
case, and this is the more consequential of the two observations.
Remark~\ref{rem:baseline} guarantees only that $\mathbf{D}=\mathbf{I}$ recovers
the control law of Chapter~\ref{ch:herpvs}; it says nothing about the case
$\mathbf{D}\neq\mathbf{I}$ with no outlier present, and the run shows that
case to be far from neutral. Every indicator moves in the same direction.
The commanded velocity at iteration~600 is $2.88\times10^{-5}$ against
$3.68\times10^{-4}\,\mathrm{m\,s^{-1}}$ for Variant~2, a factor of thirteen;
the adaptive gain has fallen further, $\lambda=0.0365$ against $0.0572$; the
intensity error is $1.91$ against $2.40$ grey levels root-mean-square; and
the rotational pose error is smaller by a factor of $6.8$. Variant~3 has not
merely stalled at a different point --- it has descended further along the
same cost function.

The mechanism, however, is not the one Section~\ref{sec:hermite-gate} claims.
%%% ---------------------------------------------------------------
%%% This paragraph is not optional. The gate as implemented derives
%%% structG from w_h/w_v/w_d, which are the image filtered by the summed
%%% Hermite kernels; those kernels have a DC gain equal to their L1 norm
%%% to three decimal places, so structG tracks local INTENSITY, not local
%%% structure. Measured on this run, normalization inflates the residual
%%% energy 15.5-fold, which against the floor of 0.1 on Gbar places of
%%% order a sixth of the image at or near the clamp. Writing the result
%%% above as a vindication of structure-aware weighting would therefore
%%% overclaim, and a reader who inspects the kernels will find it out.
%%% Either re-run with the DC-free first-order kernels and report that, or
%%% keep this paragraph. Do not delete it and keep the result.
%%% ---------------------------------------------------------------
The summed Hermite kernels $\mathbf{D}_{nm}$ are low-pass dominated --- their
DC gain equals their $L_1$ norm to three decimal places --- so the quantity
$G$ of~\eqref{eq:structG} tracks local \emph{intensity} rather than local
structure, and the normalization amplifies residuals in dark regions rather
than in geometrically flat ones. On this image the normalization inflates the
residual energy $15.5$-fold, which given the floor of $0.1$ imposed
on~\eqref{eq:gbar} places of the order of a sixth of the image at the clamp.
What the nominal run establishes, then, is that a per-pixel heteroscedastic
rescaling of the residual improves the convergence of the robust scheme
substantially; it does not yet establish that the improvement comes from
structure. Section~\ref{sec:dc-free} isolates that question by repeating the
experiment with a first-order, and therefore exactly DC-free, structure
measure.

The comparison with the unweighted baseline separates cleanly into the two
degrees of freedom. In translation, Variant~1 is the most accurate of the
three, $0.34$~mm against $0.81$ and $0.83$~mm, the deficit of both weighted
variants lying almost entirely in the depth component. This is the cost
anticipated in Section~\ref{sec:tukey-limitation}: while the pose error is
still large, valid measurements in photometrically sensitive regions produce
legitimately large residuals, a fraction of which a redescending estimator
rejects, so the weighted variants discard information the unweighted scheme
retains --- and depth, being the direction the luminance interaction matrix
constrains least, is the first to suffer. In rotation the ordering reverses:
Variant~1's $0.22^{\circ}$ is bettered by Variant~3's $0.089^{\circ}$, so the
normalized weighting does not merely recover the accuracy that plain Tukey
weighting costs but improves on the baseline. Taken together, the nominal case
shows that robust weighting is not free in the absence of outliers, as
Section~\ref{sec:tukey-limitation} predicted, but that the penalty is confined
to depth once the residual is normalized.

\begin{table}[t]
    \centering
    \caption{Nominal case (unoccluded target): final pose error, from the
    desired pose $(0,0,1300)\,\mathrm{mm}$, $(0,0,0)^{\circ}$. Translations in
    millimetres, rotations in degrees. No run reached the absolute threshold on
    $\norm{\mathbf{e}}$; all three satisfied
    $\norm{\mathbf{V}_c}<10^{-3}\,\mathrm{m\,s^{-1}}$. All values are taken at
    iteration~600, images of $320\times240$ pixels.}
    \label{tab:ch4-nominal}
    \small
    \begin{tabular}{@{}lrrrrrrrr@{}}
        \toprule
        & \multicolumn{3}{c}{Translation (mm)}
        & \multicolumn{3}{c}{Rotation (deg)}
        & \multicolumn{2}{c}{Norm} \\
        \cmidrule(lr){2-4}\cmidrule(lr){5-7}\cmidrule(lr){8-9}
        Variant & $\Delta T_x$ & $\Delta T_y$ & $\Delta T_z$
                & $\Delta\omega_x$ & $\Delta\omega_y$ & $\Delta\omega_z$
                & $\norm{\Delta\mathbf{t}}$ & $\norm{\Delta\boldsymbol{\omega}}$ \\
        \midrule
        1 --- none    & $+0.053$ & $-0.140$ & $-0.302$ & $+0.168$ & $-0.139$ & $+0.015$ & $0.337$ & $0.219$ \\
        2 --- Tukey   & $+0.004$ & $-0.166$ & $-0.790$ & $+0.568$ & $-0.196$ & $-0.007$ & $0.807$ & $0.601$ \\
        3 --- Hermite & $+0.098$ & $+0.136$ & $+0.808$ & $-0.072$ & $+0.048$ & $+0.022$ & $0.825$ & $\mathbf{0.089}$ \\
        \bottomrule
    \end{tabular}
\end{table}

%%% OPEN: whether Variants 2 and 3 genuinely coincide. The two consoles
%%% supplied agree to every printed digit on |e|, lambda, mu, |Vc| and all
%%% six pose components. Three explanations remain, in decreasing order of
%%% likelihood:
%%%   (1) the dimension guard in computeHermiteNormalizedResidual() fires,
%%%       so Variant 3 runs as Variant 2 -- test with diag_v3_gate.cpp;
%%%   (2) the two consoles are two runs of the same binary;
%%%   (3) the normalization genuinely cancels in the MAD, which is the only
%%%       one of the three that would support the prose above.
%%% Settle (1) before anything else: it is a one-line test.
